\DSource{0134}{45}
\hypertarget{AB01-PDF0134-P01}{}
\DLine{AB01-PDF0134-body-L001}{}{}{describamus, ut facile possit aequatio motus Solis inveniri. Ostendit Ptolemaeus (1) motus va-}
\DLine{AB01-PDF0134-body-L002}{}{p. 68.}{rios duobus modis effingi posse; quorum alter est ut ponatur planetam habere circulum eclip-}
\DLine{AB01-PDF0134-body-L003}{}{}{ticae concentricum, super quem alius circulus sit, cuius centrum circumferentiam primi per-}
\DLine{AB01-PDF0134-body-L004}{}{}{currat. Hic secundus circulus est circulus minor, Terram non ambiens, cuius centrum circulus}
\DLine{AB01-PDF0134-body-L005}{5}{}{maior rotare facit versus partem successionis signorum iuxta quantitatem motus longitudinalis}
\DLine{AB01-PDF0134-body-L006}{}{}{planetae versus partem in quam signa succedunt. Sidus ipsum in epicyclo, qui est circulus}
\DLine{AB01-PDF0134-body-L007}{}{}{minor, versus praecedentem aut versus subsequentem partem movetur; aut circulus minor}
\DLine{AB01-PDF0134-body-L008}{}{}{planetam versus alterutram partem rotare facit. Hic est motus anomaliae quae ad planetam}
\DLine{AB01-PDF0134-body-L009}{}{}{pertinet. --- Alter explicandi modus est ut ponatur sidus habere circulum concentricum ecli-}
\DLine{AB01-PDF0134-body-L010}{10}{}{pticae, et alterum excentricum sed illi aequalem, qui eum duobus locis secet Planeta est in}
\DLine{AB01-PDF0134-body-L011}{}{}{excentrico, sive hic eum movet, sive ipse in excentrico movetur. --- Utraque ratio anomaliam}
\DLine{AB01-PDF0134-body-L012}{}{}{pariter explicat; nos a prima incipiemus.}
\hypertarget{AB01-PDF0134-P02}{}
\DLine{AB01-PDF0134-body-L013}{}{}{\Indent Circulum ABCD eclipticae circa centrum E describamus; sit A, initio, centrum epicycli}
\DLine{AB01-PDF0134-body-L014}{}{}{HF (2). Diametrum AC usque ad H producamus, quod est punctum apogei in epicyclo; sit F}
\DLine{AB01-PDF0134-body-L015}{15}{}{locus Solis in epicyclo, et ab eo perpendicularem supra lineam AH usque ad punctum M du-}
\DLine{AB01-PDF0134-body-L016}{}{}{camus; denique lineam AF signemus, lineae AH aequalem, nam ambae sunt semidiametri}
\DLine{AB01-PDF0134-body-L017}{}{}{epicycli. Ob ea quae in hoc ipso capite iam diximus, patet eam $2^p4'\,\qfrac{3}{4}$ complecti, ut semi-}
\DLine{AB01-PDF0134-body-L018}{}{}{diametrum EF epicycli in figura praecedenti. Quo cognito, motum Solis in epicyclo versus}
\DLine{AB01-PDF0134-body-L019}{}{}{partem contrariam successioni signorum observa, vel quantitatem qua epicyclus motu medio}
\DLine{AB01-PDF0134-body-L020}{20}{}{diurno Solem versus illam partem movet, iuxta rationem qua circumferentia epicycli in $360^\circ$}
\DLine{AB01-PDF0134-body-L021}{}{}{dividitur. Motus medius Solis, qui observationibus invenitur, est motus centri epicycli versus}
\DLine{AB01-PDF0134-body-L022}{}{p. 69.}{partem successionis signorum; et hic motus quoque ea ratione ponitur qua circulus ABCD}
\DLine{AB01-PDF0134-body-L023}{}{}{$360^\circ$ complectitur.}
\hypertarget{AB01-PDF0134-P03}{}
\DLine{AB01-PDF0134-body-L024}{}{}{\Indent Arcus HF, qui est arcus epicycli inter locum Solis et punctum apogei, $30^\circ$ contineat e}
\DLine{AB01-PDF0134-body-L025}{25}{}{gradibus quorum epicyclus 360 continet. Lineam EF in hac figura ducamus; et in arcum}
\DLine{AB01-PDF0134-body-L026}{}{}{lineae FM, qui est variatio motus Solis eo loco, inquiramus. Linea EA, semidiametrus circuli}
\DLine{AB01-PDF0134-body-L027}{}{}{concentrici eclipticae, 60 partes habet ex iis quarum diametrus AC 120 continet; linea igitur}
\DLine{AB01-PDF0134-body-L028}{}{}{EH, quae centrum circuli parecliptici (3) cum apogeo epicycli coniungit, $62^p4'45''$ erit. Et}
\DLine{AB01-PDF0134-body-L029}{}{}{quia triangulum FMA est rectangulum, quadratum lineae AF erit ut summa quadratorum}
\DLine{AB01-PDF0134-body-L030}{30}{}{AM, FM; angulus autem MAF (4) notus est, ergo etiam linea FM est nota, a qua deinde co-}
\par\vspace{6pt}\hrule height.25pt\vspace{5pt}
\noindent\begin{minipage}[t]{.48\textwidth}\vspace{0pt}\fontsize{9}{11.5}\selectfont\setlength{\GutterOffset}{0pt}
\hypertarget{AB01-PDF0134-N01}{}
\DLine{AB01-PDF0134-notes_left-L001}{}{}{\Indent (1) {\itshape Almag.} III, 3, 4 et 5 (ed. Halma, t. I, p. 170-}
\DLine{AB01-PDF0134-notes_left-L002}{}{}{-183, 183--190, 190--199); quoad Solem tamen hypo-}
\DLine{AB01-PDF0134-notes_left-L003}{}{}{thesim excentrici commodiorem censet.}
\hypertarget{AB01-PDF0134-N02}{}
\DLine{AB01-PDF0134-notes_left-L004}{}{}{\Indent (2) Cod. et Plato HFT; in eorum figura enim}
\DLine{AB01-PDF0134-notes_left-L005}{}{}{punctum T ponitur loco intersectionis epicycli et}
\DLine{AB01-PDF0134-notes_left-L006}{}{}{circuli magni, inter A et D. Sed ex iis quae sequun-}
\DLine{AB01-PDF0134-notes_left-L007}{}{}{tur patet T esse punctum circuli magni indicans di-}
\DLine{AB01-PDF0134-notes_left-L008}{}{}{rectionem apogei.}
\hypertarget{AB01-PDF0134-N03}{}
\DLine{AB01-PDF0134-notes_left-L009}{}{}{\Indent (3) Scilicet paralleli seu concentrici eclipticae.}
\DLine{AB01-PDF0134-notes_left-L010}{}{}{Cum nullum habeat nomen apud scriptores nostra-}
\DLine{AB01-PDF0134-notes_left-L011}{}{}{tes, ita verto Arabicum {\itshape al-mumaththal bi falak}}
\DLine{AB01-PDF0134-notes_left-L012}{}{}{{\itshape al-burūǵ} « similem factum eclipticae » seu {\itshape al-mu-}}
\DLine{AB01-PDF0134-notes_left-L013}{}{}{{\itshape maththal} tantum, qui apud Ptolemaeum interdum}
\DLine{AB01-PDF0134-notes_left-L014}{}{}{\textgreek{ὁ ὁμόκεντρος τῷ κόσμῳ κύκλος} « circulus mundo con-}
\DLine{AB01-PDF0134-notes_left-L015}{}{}{« centricus » vel \textgreek{ὁ ὁμόκεντρος τῷ διὰ μέσων τῶν ζω-}}
\DLine{AB01-PDF0134-notes_left-L016}{}{}{\textgreek{δίων κύκλος} « circulus concentricus circulo per me-}
\end{minipage}
\hfill\begin{minipage}[t]{.48\textwidth}\vspace{0pt}\fontsize{9}{11.5}\selectfont\setlength{\GutterOffset}{0pt}
\DLine{AB01-PDF0134-notes_right-L001}{}{}{dium signorum transeunti » (i. e. eclipticae) vocatur.}
\DLine{AB01-PDF0134-notes_right-L002}{}{}{In astronomia Syriaca Barhebraei vertenda, \Name{Nau}}
\DLine{AB01-PDF0134-notes_right-L003}{}{}{eum {\itshape intersphère de la similitude} appellat; \Name{Golius},}
\DLine{AB01-PDF0134-notes_right-L004}{}{}{in versione al-Farghānī, et \Name{Sédillot}, in versione}
\DLine{AB01-PDF0134-notes_right-L005}{}{}{Ulugh Beg, {\itshape sphaeram} vel {\itshape circulum homocentri-}}
\DLine{AB01-PDF0134-notes_right-L006}{}{}{{\itshape cum}. --- Est circulus concentricus eclipticae in cuius}
\DLine{AB01-PDF0134-notes_right-L007}{}{}{plano iacet; descriptus est in superficie sphaerae}
\DLine{AB01-PDF0134-notes_right-L008}{}{}{Solis vel planetarum, {\itshape al-mumaththal} quoque vo-}
\DLine{AB01-PDF0134-notes_right-L009}{}{}{cata, quae est corpus solidum, concentricum eclip-}
\DLine{AB01-PDF0134-notes_right-L010}{}{}{ticae, duabus superficiebus parallelis contentum, qua-}
\DLine{AB01-PDF0134-notes_right-L011}{}{}{rum exterior superficiem exteriorem sphaerae circuli}
\DLine{AB01-PDF0134-notes_right-L012}{}{}{excentrici (deferentis) in puncto apogeo excentrici}
\DLine{AB01-PDF0134-notes_right-L013}{}{}{tangit; interior superficiem interiorem sphaerae ex-}
\DLine{AB01-PDF0134-notes_right-L014}{}{}{centrici in perigeo. --- Movetur ut signa zodiaci ob}
\DLine{AB01-PDF0134-notes_right-L015}{}{}{praecessionem moventur, ab ortu ad occasum.}
\hypertarget{AB01-PDF0134-N04}{}
\DLine{AB01-PDF0134-notes_right-L016}{}{}{\Indent (4) Cod. HEF.}
\end{minipage}
